\subsection{信号预处理}
\label{sec:preprocessing}

原始激光超声信号在馈入网络前经过标准化预处理流程，采集时采用有限的次平均（$N = 1$、$2$、$4$或$8$）。在任何阶段均不使用高次平均信号作为输入。

\textbf{带通滤波。}应用有限脉冲响应（FIR）带通滤波器以保留感兴趣的兰姆波频带。通带设置为$100\,\text{kHz}$--$2\,\text{MHz}$，阻带衰减为$-40\,\text{dB}$，通带纹波低于$0.1\,\text{dB}$。该范围捕获了在板试样中通常激发的基本对称（$S_0$）和反对称（$A_0$）兰姆波模式，同时抑制低频机械振动和高频电磁噪声。

\textbf{时窗选通。}矩形窗选择包含兰姆波波包的到达区域。该门在基于$S_0$模式速度（铝板$c_S \approx 5400\,\text{m/s}$）的理论上首次到达前$2\,\mu\text{s}$开始，延伸$30\,\mu\text{s}$以捕获完整波包，包括可能的多模态到达。该窗外的信号被置零以消除外部干扰。

\textbf{幅值归一化。}逐次$z$分数归一化将信号幅值标准化：
\begin{equation}
x_{\text{norm}} = \frac{x - \mu_x}{\sigma_x},
\end{equation}
其中$\mu_x$和$\sigma_x$为在整个时窗段落计算的平均值和标准差。这确保网络接收的输入具有零均值和单位方差，无论跨次采集的激光能量波动如何。

\textbf{时频变换。}时频分支在 STFT 表示上运作。STFT 参数为：汉宁窗长度$W = 256$样本；跳跃大小$H = 64$样本（75\% 重叠）；频率分辨率$\Delta f = f_s / W$，其中$f_s$为采样率。这些设置在时间定位（约$0.8\,\mu\text{s}$）和频率分辨率（约$60\,\text{kHz}$）之间取得平衡，足以分辨色散兰姆波分量。

\textbf{重采样。}所有时窗信号通过线性插值重采样至固定长度$L = 1024$样本。这为所有数据集强制执行一致的输入维度，由 PMDF-Net 的固定尺寸架构要求。